% 5. Открытый вопрос
\newcommand{\qcard}[3]{%
  \begin{navycard}[height=1.25cm,valign=center]
    \iconitem{accent}{#1}{{\bfseries\color{white}#2}\par\vspace{0.4mm}{\scriptsize\color{cardline}#3}}
  \end{navycard}}
{\begin{darkslide}
\begin{frame}
\vspace{12mm}
{\scriptsize\bfseries\color{accentlt}Открытый вопрос}\par\vspace{3mm}
{\headfont\bfseries\fontsize{17}{21}\selectfont\color{white}Возможна ли атака уклонения по\\«жёсткой метке» в чистом чёрном ящике?}\par\vspace{7mm}
\begin{columns}[T,onlytextwidth]
\column{0.485\textwidth}
  \vspace{2mm}\qcard{ic_random_w.png}{Протокол-агностична}{любой трафик: шифрованный, туннель, обычный}\par\vspace{5mm}
  \qcard{ic_eye_w.png}{Только «прошёл / не прошёл»}{единственный сигнал — блокирует ли система пакеты}
\column{0.485\textwidth}
  \vspace{2mm}\qcard{ic_puzzle_w.png}{Задаче-агностична}{сканирование, DoS, ботнеты, веб-атаки}\par\vspace{5mm}
  \qcard{ic_check_w.png}{Сохранение вредоносности}{минимальные правки, атака продолжает работать}
\end{columns}
\end{frame}
\end{darkslide}}
